\section{Preliminaries}
\label{sec:prelim}

Throughout, $H$ denotes a finite connected target graph, possibly with
loops, on $m = |V(H)|$ vertices, with adjacency matrix $A = A(H)$; a
loop at $u$ contributes $A_{uu} = 1$. Source graphs are trees on $n$
vertices. All logarithms are natural and all matrices are real.

\subsection{Walks are the only thing being counted}

Let $\mathbf{1}$ be the all-ones vector and set

\[
  w_k = A^k \mathbf{1},
  \qquad\text{so that}\qquad
  w_k(u) = \bigl|\{\text{walks of length } k \text{ in } H \text{ starting at } u\}\bigr|.
\]
A homomorphism from a tree is determined by choosing an image for a root
and then extending independently along each branch, so every count we
need is a sum of products of entries of the $w_k$.

\begin{lemma}[Path and spider counts]
  \label{lem:transfer}
  Let $\Sp(a_1,\dots,a_r)$ denote the \emph{spider}: the tree obtained by
  identifying one endpoint of $r$ paths with $a_1, \dots, a_r$ edges
  respectively at a common centre. A spider is a tree with one branch
  vertex; we write $\Sp(a,b,c)$ for the spider with leg lengths $a,b,c$.
  Root the spider at its centre. For each vertex $u \in V(H)$, count the
  homomorphisms that send the root to $u$; the legs then extend
  independently, so we sum over all such $u$.
  Then
  \[
    \Hom(P_n, H) \;=\; \mathbf{1}^{T} A^{\,n-1} \mathbf{1}
      \;=\; \sum_{u \in V(H)} w_{n-1}(u),
    \qquad
    \Hom\bigl(\Sp(a_1,\dots,a_r), H\bigr)
      \;=\; \sum_{u \in V(H)} \prod_{i=1}^{r} w_{a_i}(u).
  \]
\end{lemma}

\begin{proof}
  Root the spider at its centre. For each vertex $u \in V(H)$, count the
  homomorphisms that send the root to $u$; the branches are vertex-disjoint
  apart from the centre, so their images may be chosen independently. The
  choices along a branch with $a_i$ edges are exactly the walks of length
  $a_i$ from $u$. The path is the case $r = 2$, or directly the $k$-step
  iteration of $A$.
\end{proof}

\begin{figure}[H]
  \centering
  \begin{tikzpicture}[line cap=round, line join=round, every node/.style={font=\small}]
    \coordinate (root) at (0,0);
    \coordinate (leftend) at (-5.2,0);
    \coordinate (rightend) at (5.2,0);
    \coordinate (downend) at (0,-2.8);
    \draw[line width=0.8pt] (root) -- (leftend);
    \draw[line width=0.8pt] (root) -- (rightend);
    \draw[line width=0.8pt] (root) -- (downend);
    \foreach \x in {-4.2,-3.2,-2.2,-1.2,1.2,2.2,3.2,4.2} {
      \fill (\x,0) circle (1.8pt);
    }
    \foreach \y in {-0.9,-1.8,-2.7} {
      \fill (0,\y) circle (1.8pt);
    }
    \fill (root) circle (2.8pt);
    \node[above] at (0,0.45) {root};
    \node[left] at (-5.45,0) {$a$};
    \node[right] at (5.45,0) {$b$};
    \node[below] at (0,-3.15) {$c$};
    \node[below] at (0,-3.85) {spider $\,\Sp(a,b,c)$};
  \end{tikzpicture}
  \caption{A spider rooted at its centre.}
  \label{fig:spider-lemma}
\end{figure}

\subsection{The near-path family}

\begin{definition}[Near-path trees]
  \label{def:nearpath}
  For $n \ge 5$ and $1 \le d \le n-3$, the \emph{near-path tree} is the
  spider
  \[
    E_n^{(d)} \;=\; \Sp\bigl(n-d-2,\; d,\; 1\bigr),
  \]
  the tree on $n$ vertices obtained from the path $P_{n-1}$ by attaching
  a pendant vertex to the vertex at distance $d$ from one end. We write
  $E_n = E_n^{(2)}$ for the classical case, the tree used by
  Leontovich~\cite{leontovich1989}, obtained by hanging a pendant on the
  third-from-last vertex of $P_{n-1}$.
\end{definition}

Combining Definition~\ref{def:nearpath} with Lemma~\ref{lem:transfer},
\begin{equation}
  \label{eq:nearpath-count}
  \Hom\bigl(E_n^{(d)}, H\bigr)
    \;=\; \sum_{u \in V(H)} w_{n-d-2}(u)\, w_{d}(u)\, w_{1}(u),
\end{equation}
which is the identity implemented by every verification script in this
work. Two consequences are worth recording.

\begin{remark}[The depth parameter is a leg length]
  \label{rem:d-symmetry}
  A spider does not remember the order of its legs, so
  $E_n^{(d)} \cong E_n^{(n-2-d)}$. The parameter $d$ is therefore not a
  depth so much as the length of the shorter of the two long legs, and a
  sweep over $d \le d_{\max}$ covers exactly those near-path trees with
  $\min\{d,\, n-2-d\} \le d_{\max}$. Our sweeps use
  $d_{\max} = 20$ with $n \le 200$; the untested trees are those whose
  pendant sits near the middle of the path.
\end{remark}

\begin{remark}[Why this family and not another]
  \label{rem:why-nearpath}
  The near-path trees are the natural first place to look, for a reason
  made precise in Remark~\ref{rem:filter}: any tree that beats the path must
  match its exponential growth rate $\lambda_1^{\,n}$ and win on the
  constant in front. Branching in the interior raises the growth rate
  outright, so a challenger must confine its defect to the boundary.
  This is a heuristic for choosing a search family, not a theorem, and
  we treat every negative result obtained with it accordingly.
\end{remark}

\subsection{Symmetric trees and their quotients}

\begin{definition}[Spherically symmetric trees]
  \label{def:symtree}
  For $d_1,\dots,d_k \ge 1$, the tree $T(d_1,\dots,d_k)$ is rooted, of
  depth $k$, in which every vertex at depth $i-1$ has exactly $d_i$
  children. Its vertex set falls into $k+1$ \emph{orbits} by depth, of
  sizes $s_0 = 1$ and $s_i = d_1 d_2 \cdots d_i$, and
  \[
    |V| \;=\; 1 + d_1 + d_1 d_2 + \cdots + d_1 d_2 \cdots d_k .
  \]
  We write $\hat T(d_1,\dots,d_k)$ for the same tree with a loop added at
  the root. The case $k = 3$ recovers the classical family
  $T(x,y,z)$ on $1 + x + xy + xyz$ vertices, whose four orbits are root,
  children, grandchildren and leaves.
\end{definition}

The loop matters: it destroys bipartiteness, and with it the parity
obstruction of Section~\ref{sec:leontovich}.

\begin{definition}[Equitable partition and quotient matrix]
  \label{def:quotient}
  A partition $\pi = \{C_1,\dots,C_K\}$ of $V(H)$ is \emph{equitable} if
  there are integers $q_{ij}$ such that every vertex of $C_i$ has exactly
  $q_{ij}$ neighbours in $C_j$. The matrix $Q = (q_{ij})$ is the
  \emph{quotient matrix} of $\pi$ and $\mathbf{s} = (|C_1|,\dots,|C_K|)$
  its \emph{size vector}. The orbit partition of any automorphism group
  is equitable; so is the coarsest partition produced by color refinement.
\end{definition}

\begin{lemma}[Reduction to the quotient]
  \label{lem:quotient}
  Let $\pi$ be equitable with quotient $Q$ and size vector $\mathbf{s}$.
  Then each walk vector $w_k$ is constant on cells, and its cell values
  are the entries of $Q^k \mathbf{1}_K$. Consequently, for every $n$ and
  every spider,
  \[
    \Hom(P_n,H) \;=\; \mathbf{s}^{T} Q^{\,n-1}\mathbf{1}_K,
    \qquad
    \Hom\bigl(\Sp(a_1,\dots,a_r),H\bigr)
      \;=\; \sum_{i=1}^{K} s_i \prod_{j=1}^{r} \bigl(Q^{a_j}\mathbf{1}_K\bigr)_i .
  \]
\end{lemma}

\begin{proof}
  Induct on $k$. The claim holds for $w_0 = \mathbf{1}$. If $w_{k}$ takes
  the value $x_i$ on $C_i$, then for $u \in C_i$,
  $w_{k+1}(u) = \sum_j q_{ij} x_j$, which depends only on $i$. Summing
  Lemma~\ref{lem:transfer} cell by cell gives the two displayed
  identities.
\end{proof}

Lemma~\ref{lem:quotient} is what makes the computations feasible: a
target on $6{,}806$ vertices becomes an $18 \times 18$ integer matrix
with no loss of information about $\Hom(T,\cdot)$ for the trees we test.
Only the eigenvalues of $Q$ can appear in these counts; the rest of
$A(H)$ is invisible here, and we call those eigenvalues \emph{inert}.

\begin{lemma}[Two-step recurrence from a symmetric active spectrum]
  \label{lem:active_recurrence}
  If the active spectrum of $Q$ is contained in
  $\{\lambda_1, 0, -\lambda_1\}$, then the quotient-walk vectors satisfy
  \[
    Q^{k+2}\mathbf{1}_K \;=\; \lambda_1^{\,2}\, Q^k\mathbf{1}_K
  \]
  for every $k \ge 1$. Equivalently, after the kernel component has
  vanished, the walk counts obey $w_{k+2} = \lambda_1^2 w_k$.
\end{lemma}

\begin{proof}
  Expand $\mathbf{1}_K$ in the $\langle\cdot,\cdot\rangle_{\mathbf{s}}$-
  orthonormal eigenbasis of $Q$. The zero-eigenspace contributes nothing
  once $k \ge 1$, and the $\pm \lambda_1$ terms pick up the same factor
  $\lambda_1^2$ after two more steps.
\end{proof}

\begin{lemma}[The quotient is self-adjoint in the size weighting]
  \label{lem:selfadjoint}
  Define $\ip{x}{y} = \sum_i s_i x_i y_i$. Then $Q$ is self-adjoint with
  respect to $\langle\cdot,\cdot\rangle_{\mathbf{s}}$; in particular its
  eigenvalues are real and it has a $\langle\cdot,\cdot\rangle_{\mathbf{s}}$-orthonormal
  eigenbasis.
\end{lemma}

\begin{proof}
  Counting edges between $C_i$ and $C_j$ in two ways gives
  $s_i q_{ij} = s_j q_{ji}$, i.e.\ $\ip{Qx}{y} = \ip{x}{Qy}$.
  Equivalently $S^{1/2} Q S^{-1/2}$ is a symmetric matrix, where
  $S = \operatorname{diag}(\mathbf{s})$.
\end{proof}

\subsection{Leading coefficients and the two Leontovich properties}

Let $\lambda_1 \ge \lambda_2 \ge \cdots \ge \lambda_K$ be the eigenvalues
of $Q$ with $\langle\cdot,\cdot\rangle_{\mathbf{s}}$-orthonormal
eigenvectors $u^{(1)},\dots,u^{(K)}$, and put
$\alpha_j = \ip{\mathbf{1}}{u^{(j)}}$. Lemma~\ref{lem:quotient} and
Lemma~\ref{lem:selfadjoint} give the exact expansions
\begin{equation}
  \label{eq:expansion}
  \Hom(P_n,H) = \sum_{j=1}^{K} \alpha_j^{2}\, \lambda_j^{\,n-1},
  \qquad
  \Hom\bigl(E_n^{(d)},H\bigr)
    = \sum_{j=1}^{K} \alpha_j \beta_j^{(d)}\, \lambda_j^{\,n-d-2},
  \qquad
  \beta_j^{(d)} = \ip{w_1 \odot w_d}{u^{(j)}},
\end{equation}
where $\odot$ is the entrywise product. Everything in this paper is a
statement about the interference between the $j = 1$ term and the rest.

\begin{definition}[Leontovich, strongly Leontovich]
  \label{def:leontovich}
  A target $H$ is \emph{Leontovich} if there is an $n$ and a tree $T$ on
  $n$ vertices with $\Hom(T,H) < \Hom(P_n,H)$: the path is not a
  minimizer. It is \emph{strongly Leontovich}~\cite{galvin2026} if
  moreover a fixed near-path family beats the path for all large $n$,
  i.e.\ $\Hom(E_n^{(d)},H) < \Hom(P_n,H)$ for all sufficiently large $n$
  of the relevant parity.
\end{definition}

If $H$ is connected and non-bipartite, $\lambda_1$ strictly dominates
$|\lambda_j|$ for $j \ge 2$, and \eqref{eq:expansion} identifies the
strong property with a single inequality between leading coefficients.

\begin{definition}[Leading-coefficient ratio]
  \label{def:rho}
  For a connected non-bipartite target $H$ with active quotient
  $(Q,\mathbf{s})$, Perron eigenvalue $\lambda_1$ and positive Perron
  eigenvector $u = u^{(1)}$, set
  \begin{equation}
    \label{eq:rho}
    \rho_d(H) \;=\;
      \lim_{n \to \infty}
      \frac{\Hom\bigl(E_n^{(d)},H\bigr)}{\Hom(P_n,H)}
    \;=\;
      \frac{\ip{w_1 \odot w_d}{u}}
           {\lambda_1^{\,d+1}\, \ip{\mathbf{1}}{u}} .
  \end{equation}
  The right-hand side is independent of the normalisation of $u$. Then
  $H$ is strongly Leontovich for depth $d$ precisely when
  $\rho_d(H) < 1$, and $\rho_d(H) > 1$ forbids it.
\end{definition}

\begin{remark}[Bipartite targets and parity]
  \label{rem:bipartite-parity}
  If $H$ is bipartite then $-\lambda_1$ is also an eigenvalue of $Q$ and
  the limit in \eqref{eq:rho} need not exist: the two dominant modes
  reinforce at one parity of $n$ and cancel at the other. This is not a
  technicality but the governing phenomenon. It is why every crossover
  in this paper occurs at odd $n$ unless the target is very large
  (Section~\ref{sec:doublecover}), and why the looped root in
  $\hat T(\cdot)$ is the standard device for removing the obstruction.
  For bipartite $H$ we therefore read \eqref{eq:rho} along a fixed
  parity class and say so explicitly.
\end{remark}

\begin{remark}[The margins involved]
  \label{rem:margins}
  To fix expectations: the smallest witnesses in this paper win by a
  relative margin of about $4.5 \times 10^{-4}$ ($H_{18}$ at $(n,d)=(15,4)$)
  down to $4.8 \times 10^{-5}$ (the perturbed cover of
  Section~\ref{sec:doublecover}). A double-precision evaluation of
  \eqref{eq:nearpath-count} accumulates relative error on the order of
  $n\,m\,\varepsilon \approx 10^{-12}$, which is comfortably below these
  margins but not below every margin that might exist. Floating point is
  therefore used only to \emph{flag} candidates; every reported witness
  is re-derived in exact integer arithmetic through
  Lemma~\ref{lem:quotient}, and every asymptotic claim is accompanied by
  the high-precision error estimate described in
  Section~\ref{sec:evidence}.
\end{remark}

\subsection{Notation summary}

\begin{table}[H]
\centering
\caption{Notation used throughout.}
\label{tab:notation}
\begin{tabular}{@{}ll@{}}
  \toprule
  Symbol & Meaning \\
  \midrule
  $\Hom(G,H)$ & number of homomorphisms $G \to H$ \\
  $m$, $n$ & order of the target $H$; order of the source tree \\
  $A$, $w_k = A^k\mathbf{1}$ & adjacency matrix of $H$; walk-count vector \\
  $P_n$ & path on $n$ vertices ($n-1$ edges) \\
  $\Sp(a_1,\dots,a_r)$ & spider with legs of $a_1,\dots,a_r$ edges \\
  $E_n^{(d)} = \Sp(n{-}d{-}2, d, 1)$ & near-path tree; $E_n = E_n^{(2)}$ \\
  $T(d_1,\dots,d_k)$, $\hat T(\cdot)$ & symmetric tree; same with a looped root \\
  $Q$, $\mathbf{s}$, $K$ & quotient matrix, cell sizes, number of cells \\
  $\lambda_1 > |\lambda_j|$, $u$ & Perron eigenvalue and eigenvector of $Q$ \\
  $\lambda_2^{*}$ & largest active subdominant eigenvalue \\
  $\rho_d(H)$ & leading-coefficient ratio, \eqref{eq:rho} \\
  $\Delta_n = \Hom(P_n,H) - \Hom(E_n^{(d)},H)$ & crossover margin; $\Delta_n > 0$ means the path loses \\
  \bottomrule
\end{tabular}
\end{table}
